% Terjemahan Bahasa Indonesia dari chapter2.tex baris 911--1132.
% Sumber: Methods of Algebra, Volume 2, commit
% 9a5803ff77dd3257484cb177f851a73770a59dd3 (CC BY 4.0).
% stable-unit-id: o014.aljabr2.chapter2.subobjects-and-isomorphism-theorems

% segment-id: o014.li2.chapter2.subobjects-and-isomorphism-theorems.h001
\section{Subobjek dan Teorema Isomorfisme}\label{sec:Abel-cat-subobjects}
\hypertarget{o014.aljabr2.chapter2.subobjects-and-isomorphism-theorems}{ }

% segment-id: o014.li2.chapter2.subobjects-and-isomorphism-theorems.p001
Pada bagian ini, tetapkan suatu kategori abelian $\mathcal{A}$. Untuk setiap
objek $X$ dalam $\mathcal{A}$, himpunan subobjeknya yang terurut parsial
dilambangkan dengan $(\mathrm{Sub}_X, \subset)$, mengikuti konvensi dalam
\S\ref{sec:subquot}.

% segment-id: o014.li2.chapter2.subobjects-and-isomorphism-theorems.p002
Hasil-hasil dalam bagian ini sudah sangat dikenal untuk kasus
$\mathcal{A}=R\dcate{Mod}$, dengan $R$ sebarang gelanggang; tetapi untuk
kategori abelian diperlukan argumen yang berbeda.

% segment-id: o014.li2.chapter2.subobjects-and-isomorphism-theorems.d001
\begin{definisi}\label{def:intersection-sum}
	\index{kuosien (quotient)}
	\index[sym1]{XoverY@$X/Y$}
	Misalkan $X$ merupakan objek dari $\mathcal{A}$.
	% segment-id: o014.li2.chapter2.subobjects-and-isomorphism-theorems.l001
	\begin{itemize}
		% segment-id: o014.li2.chapter2.subobjects-and-isomorphism-theorems.i001
		\item Untuk suatu subobjek $X'$ dari $X$, \emph{kuosien} yang
		bersesuaian didefinisikan sebagai
		$X/X':=\Coker[X'\to X]$; objek ini secara alami merupakan kuosien
		dari $X$. Setiap epimorfisme
		$f:X\twoheadrightarrow X''$ dapat dipandang sebagai kuosien $X$ oleh
		$\Ker(f)$.
		% segment-id: o014.li2.chapter2.subobjects-and-isomorphism-theorems.i002
		\item Untuk keluarga subobjek $X_1,\ldots,X_n$ dari $X$, \emph{irisan}
		dan \emph{jumlah subobjek} yang bersesuaian masing-masing didefinisikan
		sebagai
		% segment-id: o014.li2.chapter2.subobjects-and-isomorphism-theorems.q001
		\begin{align*}
			X_1 \cap \cdots \cap X_n & := X_1 \dtimes{X} \cdots \dtimes{X} X_n, \\
			X_1 + \cdots + X_n & := \Image\left[ X_1 \oplus \cdots \oplus X_n \xrightarrow{\sigma} X \right].
		\end{align*}
		Di sini, $\sigma$ diinduksi oleh semua $X_i\hookrightarrow X$ (dan dapat
		dipandang sebagai penjumlahan). Keduanya secara alami merupakan subobjek
		dari $X$: kasus $X_1\cap\cdots\cap X_n$ mengikuti
		Lema~\ref{prop:mono-product}, sedangkan kasus $X_1+\cdots+X_n$ mengikuti
		langsung dari definisi.
	\end{itemize}
\end{definisi}

% segment-id: o014.li2.chapter2.subobjects-and-isomorphism-theorems.p003
Sebagai contoh, definisi kohomologi dapat ditulis ulang menggunakan kuosien:
% segment-id: o014.li2.chapter2.subobjects-and-isomorphism-theorems.q002
\[ \Hm^n(X) := \Ker(d_X^n)/\Image(d_X^{n-1}). \]

% segment-id: o014.li2.chapter2.subobjects-and-isomorphism-theorems.p004
Secara lebih umum, untuk sebarang himpunan $I$ dan keluarga subobjek
$(X_i)_{i\in I}$ dari $X$, metode yang sama memberikan definisi yang baik
bagi $\bigcap_{i\in I}X_i$ (atau $\sum_{i\in I}X_i$), asalkan produk seratnya
di atas $X$ (atau koproduknya $\bigsqcup_{i\in I}X_i$ dalam $\mathcal{A}$)
ada. Arti penting definisi ini tampak dalam fakta berikut.

% segment-id: o014.li2.chapter2.subobjects-and-isomorphism-theorems.pr001
\begin{proposisi}\label{prop:subobject-inf-sup}
	Ambil suatu keluarga subobjek $(X_i)_{i\in I}$ dari $X$. Jika produk
	seratnya di atas $X$ (atau koproduknya $\bigsqcup$ dalam $\mathcal{A}$)
	ada, masing-masing objek tersebut memberikan infimum (atau supremum) dari
	$(X_i)_{i\in I}$ dalam himpunan terurut parsial
	$(\mathrm{Sub}_X,\subset)$.
\end{proposisi}
% segment-id: o014.li2.chapter2.subobjects-and-isomorphism-theorems.pf001
\begin{bukti}
	Dari konstruksinya, untuk setiap $i$ berlaku
	$\bigcap_{j\in I}X_j\subset X_i\subset\sum_{j\in I}X_j$.

	Berikutnya, tinjau suatu subobjek $Y\hookrightarrow X$. Untuk kasus irisan,
	andaikan $\forall i\in I,\;Y\subset X_i$, yakni terdapat keluarga diagram
	komutatif
	% segment-id: o014.li2.chapter2.subobjects-and-isomorphism-theorems.g001
	$\begin{tikzcd}[row sep=small, column sep=small]
		Y \arrow[hookrightarrow, r] \arrow[d] & X \\
		X_i \arrow[hookrightarrow, ru] &
	\end{tikzcd}$.
	Sifat universal produk serat kemudian memberikan morfisme
	$Y\to\bigcap_{i\in I}X_i$, sehingga
	$Y\subset\bigcap_{i\in I}X_i$. Untuk kasus jumlah subobjek, andaikan
	$\forall i\in I,\;X_i\subset Y$. Sifat universal koproduk menghasilkan
	diagram komutatif
	% segment-id: o014.li2.chapter2.subobjects-and-isomorphism-theorems.g002
	\[\begin{tikzcd}
		\coprod_{i \in I} X_i \arrow[r, "\sigma"] \arrow[rd] & X \\
		& Y \arrow[hookrightarrow, u].
	\end{tikzcd}\]

	Dengan menerapkan Lema~\ref{prop:Im-Coim-fac}, morfisme
	$\bigsqcup_{i\in I}X_i\to Y$ berfaktor secara unik melalui
	$\Coim(\sigma)\simeq\Image(\sigma)$, serasi dengan morfisme menuju $X$.
	Dengan demikian,
	$\sum_{i\in I}X_i:=\Image(\sigma)\subset Y$.
\end{bukti}

% segment-id: o014.li2.chapter2.subobjects-and-isomorphism-theorems.c001
\begin{konvensi}\label{con:increasing-union}
	\index[sym1]{Xsum@$\sum_{i \in I} X_i$}
	\index[sym1]{Xunion@$\bigcup_{i \in I} X_i$}
	Berdasarkan Proposisi~\ref{prop:subobject-inf-sup}, kita juga dapat
	menghindari penggunaan produk serat atau koproduk dan langsung
	mendefinisikan irisan $\bigcap_iX_i$ (atau jumlah subobjek
	$\sum_iX_i$) dari suatu keluarga subobjek $(X_i)_{i\in I}$ sebagai
	infimum (atau supremum) keluarga tersebut dalam $\mathrm{Sub}_X$, asalkan
	infimum (atau supremum) itu ada. Di sini selalu diandaikan bahwa himpunan
	indeks $I$ merupakan himpunan kecil. Jika $I$ dilengkapi dengan urutan
	parsial terfilter dan $i\leq j\implies X_i\subset X_j$, supremum
	$\sum_{i\in I}X_i$ dapat pula dilambangkan secara wajar sebagai
	gabungan naik $\bigcup_{i\in I}X_i$.
\end{konvensi}

% segment-id: o014.li2.chapter2.subobjects-and-isomorphism-theorems.co001
\begin{korolari}
	Untuk setiap objek $X$, himpunan terurut parsial
	$(\mathrm{Sub}_X,\subset)$ merupakan kisi terbatas dalam pengertian
	Definisi~\ref{def:lattice-order}.
\end{korolari}
% segment-id: o014.li2.chapter2.subobjects-and-isomorphism-theorems.pf002
\begin{bukti}
	Supremum dari dua subobjek $Y,Z$ adalah $Y+Z$, sedangkan infimumnya adalah
	$Y\cap Z$; unsur terbesar dalam $\mathrm{Sub}_X$ adalah $X$, dan unsur
	terkecilnya adalah $0$.
\end{bukti}

% segment-id: o014.li2.chapter2.subobjects-and-isomorphism-theorems.p005
Bagian ini berfokus pada struktur $(\mathrm{Sub}_X,\subset)$, tetapi versi
dalam hal objek kuosien $(\mathrm{Quot}_X,\twoheadleftarrow)$ sepenuhnya
sepadan. Hal ini mengikuti dari bijeksi pembalik urutan yang jelas
$\mathrm{Sub}_X\simeq\mathrm{Quot}_X$: subobjek $X'$ dan objek kuosien $X''$
saling bersesuaian jika dan hanya jika keduanya dapat ditempatkan dalam suatu
barisan eksak pendek $0\to X'\to X\to X''\to0$. Karena itu, versi tersebut
tidak akan dibahas secara terpisah.

% segment-id: o014.li2.chapter2.subobjects-and-isomorphism-theorems.d002
\begin{definisi}\label{def:subobject-image}
	\index[sym1]{finvY@$f^{-1}(Y')$}
	\index[sym1]{fX@$f(X')$}
	Diberikan morfisme $f:X\to Y$ serta subobjek-subobjek
	$X'\hookrightarrow X$ dan $Y'\hookrightarrow Y$, tuliskan
	% segment-id: o014.li2.chapter2.subobjects-and-isomorphism-theorems.q003
	\[ f^{-1}(Y') := X \dtimes{Y} Y', \quad f(X') := \Image\left[ X' \hookrightarrow X \xrightarrow{f} Y \right]; \]
	objek pertama merupakan prabayang $Y'$ dan subobjek dari $X$, sedangkan
	objek kedua merupakan citra $X'$ dan subobjek dari $Y$. Konstruksi ini
	memberikan pemetaan yang mempertahankan urutan dalam dua arah:
	% segment-id: o014.li2.chapter2.subobjects-and-isomorphism-theorems.g003
	$\begin{tikzcd}
		\mathrm{Sub}_X \arrow[r, yshift=0.3em, "f(\cdot)"] & \mathrm{Sub}_Y \arrow[l, yshift=-0.3em, "f^{-1}(\cdot)"]
	\end{tikzcd}$.
\end{definisi}

% segment-id: o014.li2.chapter2.subobjects-and-isomorphism-theorems.le001
\begin{lema}\label{prop:image-sum-preimage-intersection}
	Diberikan $f:X\to Y$ seperti di atas.
	% segment-id: o014.li2.chapter2.subobjects-and-isomorphism-theorems.l002
	\begin{enumerate}[(i)]
		% segment-id: o014.li2.chapter2.subobjects-and-isomorphism-theorems.i003
		\item Untuk setiap $X'\in\mathrm{Sub}_X$ dan
		$Y'\in\mathrm{Sub}_Y$, berlaku
		% segment-id: o014.li2.chapter2.subobjects-and-isomorphism-theorems.q004
		\[ X' \subset f^{-1}(Y') \iff f(X') \subset Y'. \]
		% segment-id: o014.li2.chapter2.subobjects-and-isomorphism-theorems.i004
		\item Untuk keluarga subobjek $(X'_i)_{i\in I}$ dari
		$\mathrm{Sub}_X$ (atau $(Y'_i)_{i\in I}$ dari
		$\mathrm{Sub}_Y$),\footnote{Catatan penerjemah (O014-C025): sumber
		memperkenalkan keluarga $(X_i)$ dan $(Y_i)$ tanpa tanda prima, sedangkan
		kedua rumus berikut memakai $X'_i$ dan $Y'_i$. Target menambahkan tanda
		prima pada deklarasi agar cocok dengan rumus.} kita mempunyai
		% segment-id: o014.li2.chapter2.subobjects-and-isomorphism-theorems.q005
		\[ f\left( \sum_{i \in I} X'_i \right) = \sum_{i \in I} f(X'_i), \quad f^{-1}\left( \bigcap_{i \in I} Y'_i \right) = \bigcap_{i \in I} f^{-1}(Y'_i), \]
		asalkan irisan dan jumlah subobjek yang dibicarakan ada.
	\end{enumerate}
\end{lema}
% segment-id: o014.li2.chapter2.subobjects-and-isomorphism-theorems.pf003
\begin{bukti}
	Untuk (i), menurut definisi, $X'\subset f^{-1}(Y')$ ekuivalen dengan
	keberadaan diagram komutatif
	% segment-id: o014.li2.chapter2.subobjects-and-isomorphism-theorems.g004
	\[\begin{tikzcd}
		X' \arrow[dashed, r, "\exists"] \arrow[hookrightarrow, d] & Y' \arrow[hookrightarrow, d] \\
		X \arrow[r, "f"'] & Y .
	\end{tikzcd}\]
	Karena $f(X')$ adalah citra dari
	$X'\hookrightarrow X\xrightarrow{f}Y$,\footnote{Catatan penerjemah
	(O014-C026): sumber menyebut $f(X')$ sebagai kocitra, tetapi
	Definisi~\ref{def:subobject-image} menetapkannya sebagai
	$\Image[X'\to X\to Y]$, yakni citra dan subobjek dari $Y$. Target memakai
	istilah yang sesuai dengan objek tersebut.} Lema~\ref{prop:Im-Coim-fac}
	menunjukkan bahwa diagram seperti itu berpadanan secara bijektif dengan
	% segment-id: o014.li2.chapter2.subobjects-and-isomorphism-theorems.g005
	\[\begin{tikzcd}
		X' \arrow[r] \arrow[d, hookrightarrow] & f(X') \arrow[hookrightarrow, d] \arrow[dashed, r, "\exists"] & Y' \arrow[hookrightarrow, ld] \\
		X \arrow[r, "f"'] & Y &
	\end{tikzcd}\]
	yang membuktikan (i). Jika himpunan-himpunan terurut parsial tersebut
	dipandang sebagai kategori $\cate{Sub}_X$ dan seterusnya, (i) menyatakan
	bahwa $f(\cdot):\cate{Sub}_X\to\cate{Sub}_Y$ merupakan funktor adjoin kiri
	bagi $f^{-1}(\cdot):\cate{Sub}_Y\to\cate{Sub}_X$. Karena produk (atau
	koproduk) bersesuaian dengan infimum (atau supremum), (ii) merupakan
	akibat dari \cite[Teorema 2.8.12]{Li1}.
\end{bukti}

% segment-id: o014.li2.chapter2.subobjects-and-isomorphism-theorems.pr002
\begin{proposisi}\label{prop:biproduct-sum-intersection}
	Tinjau subobjek-subobjek $i:A\hookrightarrow X$ dan
	$j:B\hookrightarrow X$. Morfisme $(i,j):A\oplus B\to X$ merupakan
	isomorfisme jika dan hanya jika
	% segment-id: o014.li2.chapter2.subobjects-and-isomorphism-theorems.q006
	\[ A \cap B = 0, \quad A + B = X. \]
	Selain itu, diagram komutatif
	% segment-id: o014.li2.chapter2.subobjects-and-isomorphism-theorems.g006
	\[\begin{tikzcd}
		A \cap B \arrow[r, "\delta_1"] \arrow[d, "\delta_2"'] & A \arrow[d, "i"] \\
		B \arrow[r, "j"'] & A + B
	\end{tikzcd}\]
	merupakan sekaligus diagram dorong keluar dan diagram tarik balik; di sini
	$\delta_1$ dan $\delta_2$ adalah morfisme kanonik.
\end{proposisi}
% segment-id: o014.li2.chapter2.subobjects-and-isomorphism-theorems.pf004
\begin{bukti}
	Definisikan morfisme diagonal $\Delta:A\cap B\to A\oplus B$ melalui
	$p_i\Delta=\delta_i$ untuk $i=1,2$. Definisikan pula morfisme
	antidiagonal $\Delta^-:A\cap B\to A\oplus B$ melalui
	$p_1\Delta^-=-\delta_1$ dan $p_2\Delta^-=\delta_2$.

	Catatan~\ref{rem:Abel-cat-fibered-product} menyatakan bahwa $\Delta$
	memberikan isomorfisme
	$A\cap B:=A\dtimes{X}B\rightiso
	\Ker\left[A\oplus B\xrightarrow{(i,-j)}X\right]$. Jika $(i,-j)$ diganti
	dengan $(i,j)$, citra $A\oplus B\to X$ tetap $A+B$, sedangkan morfisme
	kernelnya berubah dari $\Delta$ menjadi $\Delta^-$. Dengan demikian,
	diperoleh diagram komutatif
	dengan baris eksak
	% segment-id: o014.li2.chapter2.subobjects-and-isomorphism-theorems.g007
	\[\begin{tikzcd}
		& & & X & \\
		0 \arrow[r] & A \cap B \arrow[r, "{\Delta^-}"'] & A \oplus B \arrow[ru, "{(i, j)}" description] \arrow[r] & A+B \arrow[r] \arrow[hookrightarrow, u] & 0
	\end{tikzcd}\]
	Oleh karena itu, $(i,j):A\oplus B\to X$ merupakan isomorfisme jika dan
	hanya jika $A+B=X$ dan $A\cap B=0$.

	Catatan~\ref{rem:Abel-cat-fibered-product} juga mengidentifikasi
	$A\dsqcup{A\cap B}B$ dengan
	$(A\oplus B)/\Image(\Delta^-)$. Jadi, kita mempunyai diagram komutatif
	% segment-id: o014.li2.chapter2.subobjects-and-isomorphism-theorems.g008
	\[\begin{tikzcd}[column sep=small]
		A \cap B \arrow[r, "\delta_1"] \arrow[d, "\delta_2"'] & A \arrow[d] \arrow[rdd, bend left, "i" description] & \\
		B \arrow[r] \arrow[rrd, bend right, "j" description] & (A \oplus B)/\Image(\Delta^-) \arrow[rd, "\simeq" description] & \\
		& & A + B
	\end{tikzcd}\]
	dan persegi kiri atasnya merupakan diagram dorong keluar; karena itu, seluruh
	bingkai luarnya juga merupakan diagram dorong keluar. Karena $\delta_1$ merupakan
	monomorfisme, Proposisi~\ref{prop:Abel-cat-pull-push} menunjukkan bahwa
	bingkai luar tersebut juga merupakan diagram tarik balik.
\end{bukti}

% segment-id: o014.li2.chapter2.subobjects-and-isomorphism-theorems.p006
Dengan kata lain, $B$ merupakan komplemen dari $A$ dalam kisi
$\mathrm{Sub}_X$ (Definisi~\ref{def:modular-lattice}) jika dan hanya jika
$(i,j):A\oplus B\rightiso X$; kondisi terakhir ini juga sering ditulis
sebagai persamaan $A\oplus B=X$. Berdasarkan fakta ini, kita dapat menggunakan
bahasa teori kisi secara rekursif untuk menentukan apakah $X$ dapat
dinyatakan sebagai jumlah langsung dari sejumlah hingga subobjek yang
diberikan, $X_1,\ldots,X_n$; hasilnya serupa dengan kasus teori modul. Untuk
keluarga subobjek tak hingga, pada umumnya kita perlu bekerja dalam kategori
Grothendieck; lihat
\ifcsname r@sec:Grothendieck-cat\endcsname\S\ref{sec:Grothendieck-cat}\else\S2.10\fi.
Pembahasan terkait diserahkan kepada latihan bab ini.

% segment-id: o014.li2.chapter2.subobjects-and-isomorphism-theorems.p007
Teori modul bertumpu pada beberapa teorema isomorfisme dasar; lihat
\cite[Proposisi 6.1.11, 6.1.12, 6.1.13]{Li1}. Berikut ini kita memperluasnya
ke kategori abelian umum $\mathcal{A}$.

% segment-id: o014.li2.chapter2.subobjects-and-isomorphism-theorems.th001
\begin{teorema}\label{prop:Abel-cat-isom-thm}
	Tetapkan suatu objek $X$.
	% segment-id: o014.li2.chapter2.subobjects-and-isomorphism-theorems.l003
	\begin{enumerate}[(i)]
		% segment-id: o014.li2.chapter2.subobjects-and-isomorphism-theorems.i005
		\item Untuk setiap morfisme $f:X\to Y$, terdapat isomorfisme kanonik
		$X/\Ker(f)\rightiso\Image(f)$ yang membuat diagram
		% segment-id: o014.li2.chapter2.subobjects-and-isomorphism-theorems.g009
		\[\begin{tikzcd}
			X \arrow[twoheadrightarrow, r] \arrow[twoheadrightarrow, d] & \Image(f) \\
			X/\Ker(f) \arrow[ru, "\sim" sloped]
		\end{tikzcd} \quad \text{komutatif}. \]
		% segment-id: o014.li2.chapter2.subobjects-and-isomorphism-theorems.i006
		\item Tetapkan subobjek $Z\hookrightarrow X$, dan lambangkan morfisme
		alami $X\twoheadrightarrow\overline{X}:=X/Z$ dengan $\pi$. Terdapat
		bijeksi yang mempertahankan urutan
		% segment-id: o014.li2.chapter2.subobjects-and-isomorphism-theorems.g010
		\[\begin{tikzcd}[row sep=tiny]
			\left\{ Y \in \mathrm{Sub}_X : Z \subset Y \right\} \arrow[leftrightarrow, r, "1:1"] & \mathrm{Sub}_{\overline{X}} \\
			Y \arrow[r, mapsto] & \pi(Y) = Y/Z =: \overline{Y} \\
			\pi^{-1}(\overline{Y}) & \overline{Y} \arrow[mapsto, l] ,
		\end{tikzcd}\]
		dengan notasi seperti dalam Definisi~\ref{def:subobject-image}. Jika
		$Y\supset Z$ seperti di atas, terdapat isomorfisme kanonik
		$X/Y\rightiso\overline{X}/\overline{Y}$ yang membuat diagram
		% segment-id: o014.li2.chapter2.subobjects-and-isomorphism-theorems.g011
		\[\begin{tikzcd}
			X \arrow[r] \arrow[twoheadrightarrow, d] & \overline{X} \arrow[twoheadrightarrow, d] \\
			X/Y \arrow[r, "\sim"'] & \overline{X}/\overline{Y}
		\end{tikzcd}\]
		komutatif. Selain itu,
		$\overline{Y_1\cap Y_2}=\overline{Y_1}\cap\overline{Y_2}$ dan
		$\overline{Y_1+Y_2}=\overline{Y_1}+\overline{Y_2}$.
		% segment-id: o014.li2.chapter2.subobjects-and-isomorphism-theorems.i007
		\item Untuk $Y,Z\in\mathrm{Sub}_X$, terdapat isomorfisme kanonik
		$Y/(Y\cap Z)\rightiso(Y+Z)/Z$ yang membuat diagram
		% segment-id: o014.li2.chapter2.subobjects-and-isomorphism-theorems.g012
		\[\begin{tikzcd}
			Y \arrow[hookrightarrow, r] \arrow[twoheadrightarrow, d] & Y+Z \arrow[twoheadrightarrow, d] \\
			Y/(Y \cap Z) \arrow[r, "\sim"'] & (Y+Z)/Z
		\end{tikzcd} \quad \text{komutatif}. \]
	\end{enumerate}
\end{teorema}
% segment-id: o014.li2.chapter2.subobjects-and-isomorphism-theorems.pf005
\begin{bukti}
	Mengingat barisan eksak pendek
	$0\to\Ker(f)\to X\to\Image(f)\to0$, pernyataan (i) hanyalah
	perumusan ulang definisi.

	Untuk (ii), setelah $Y$ diberikan, mula-mula bentuk diagram komutatif
	dengan baris-baris eksak
	% segment-id: o014.li2.chapter2.subobjects-and-isomorphism-theorems.g013
	\[\begin{tikzcd}
		0 \arrow[r] & Z \arrow[r] \arrow[equal, d] & Y \arrow[r] \arrow[hookrightarrow, d] & Y/Z \arrow[dashed, d, "\exists!"] \arrow[r] & 0 \\
		0 \arrow[r] & Z \arrow[r] & X \arrow[r, "\pi"'] & \overline{X} \arrow[r] & 0
	\end{tikzcd}\]
	Anak panah bergaris putus-putus berasal dari funktorialitas kokernel
	\eqref{eqn:Ker-Coker-functoriality}. Dengan menerapkan
	Teorema~\ref{prop:snake-lemma}, segera diperoleh
	$\Ker[Y/Z\to\overline{X}]=0$ dan
	$X/Y\rightiso\Coker[Y/Z\to\overline{X}]$. Khususnya,
	% segment-id: o014.li2.chapter2.subobjects-and-isomorphism-theorems.q007
	\[ Y/Z =  \Image[Y \to \overline{X}] = \pi(Y) \;\in \mathrm{Sub}_{\overline{X}}, \quad X/Y \simeq \overline{X}/\pi(Y). \]

	Sebelumnya telah dijelaskan bahwa kedua pemetaan pada (ii) mempertahankan
	urutan. Berikutnya akan dibuktikan bahwa keduanya saling invers. Untuk suatu
	$Y$, karena $\overline{Y}\hookrightarrow\overline{X}$ merupakan kernel dari
	$\overline{X}\to\overline{X}/\overline{Y}$,
	Proposisi~\ref{prop:Ker-Coker-composite} menyatakan bahwa
	$\pi^{-1}(\overline{Y}):=X\dtimes{\overline{X}}\overline{Y}\hookrightarrow X$
	merupakan kernel dari morfisme komposit
	$X\to\overline{X}\to\overline{X}/\overline{Y}$. Berkat isomorfisme di
	atas, subobjek ini juga merupakan kernel dari $X\to X/Y$. Hal ini
	mengidentifikasi $\pi^{-1}(\overline{Y})$ dengan $Y$.

	Sebaliknya, diberikan $\overline{Y}$, tetapkan
	$Y:=X\dtimes{\overline{X}}\overline{Y}=\pi^{-1}(\overline{Y})$. Tinjau
	diagram tarik balik
	% segment-id: o014.li2.chapter2.subobjects-and-isomorphism-theorems.g014
	\[\begin{tikzcd}
		Y \arrow[r] \arrow[hookrightarrow, d] & \overline{Y} \arrow[hookrightarrow, d] \\
		X \arrow[r, "\pi"'] & \overline{X} \arrow[phantom, lu, "\Box" description] .
	\end{tikzcd}\]
	Karena $\pi$ pada baris bawah merupakan epimorfisme,
	Proposisi~\ref{prop:Abel-cat-pull-push} menyatakan bahwa morfisme pada baris
	atas juga merupakan epimorfisme. Ini menunjukkan
	$\overline{Y}=\pi(Y)$.

	Terakhir, persamaan
	$\overline{Y_1\cap Y_2}=\overline{Y_1}\cap\overline{Y_2}$ dan
	$\overline{Y_1+Y_2}=\overline{Y_1}+\overline{Y_2}$ merupakan akibat
	langsung dari bijeksi yang mempertahankan urutan
	$Y\leftrightarrow\overline{Y}$.

	Untuk (iii), tinjau diagram komutatif dengan baris-baris eksak
	% segment-id: o014.li2.chapter2.subobjects-and-isomorphism-theorems.g015
	\[\begin{tikzcd}
		0 \arrow[r] & Y \cap Z \arrow[r] \arrow[hookrightarrow, d] & Y \arrow[hookrightarrow, d] \arrow[r] & Y/(Y \cap Z) \arrow[r] \arrow[dashed, d, "\exists! \theta"] & 0 \\
		0 \arrow[r] & Z \arrow[r] & Y+Z \arrow[r] & (Y+Z)/Z \arrow[r] & 0
	\end{tikzcd}\]
	dengan $\theta$ berasal dari funktorialitas kokernel.
	Proposisi~\ref{prop:biproduct-sum-intersection} menyatakan bahwa persegi
	kirinya merupakan diagram dorong keluar. Karena diagram dorong keluar mempertahankan
	kokernel (Proposisi~\ref{prop:pull-back-ker}), $\theta$ merupakan
	isomorfisme, sebagaimana dikehendaki.
\end{bukti}

% segment-id: o014.li2.chapter2.subobjects-and-isomorphism-theorems.co002
\begin{korolari}\label{prop:Abel-cat-Coker-pull}
	Tinjau diagram komutatif
	% segment-id: o014.li2.chapter2.subobjects-and-isomorphism-theorems.g016
	\[\begin{tikzcd}
		\Ker(g') \arrow[hookrightarrow, r] \arrow[d, "k"'] & W \arrow[r, "{g'}"] \arrow[d, "{f'}"'] & Y \arrow[d, "f"] \arrow[twoheadrightarrow, r] & \Coker(g') \arrow[d, "c"] \\
		\Ker(g) \arrow[hookrightarrow, r] & X \arrow[r, "g"'] & Z \arrow[twoheadrightarrow, r] & \Coker(g)
	\end{tikzcd}\]
	dengan $k,c$ merupakan morfisme yang dicirikan oleh funktorialitas
	\eqref{eqn:Ker-Coker-functoriality}. Jika persegi tengah merupakan diagram
	tarik balik (atau dorong keluar), maka $c$ merupakan monomorfisme (atau $k$
	merupakan epimorfisme).
\end{korolari}
% segment-id: o014.li2.chapter2.subobjects-and-isomorphism-theorems.pf006
\begin{bukti}
	Cukup menangani kasus diagram tarik balik. Karena penarikan balik sepanjang
	$f$ dapat dilakukan secara bertahap, cukup ditinjau dua kasus: $f$
	merupakan epimorfisme atau $f$ merupakan monomorfisme. Jika $f$ merupakan
	epimorfisme, Proposisi~\ref{prop:Abel-cat-pull-push} menyatakan bahwa
	persegi tengah juga merupakan diagram dorong keluar; dalam hal ini $c$
	merupakan isomorfisme. Jika $f$ merupakan monomorfisme, uraikan penarikan balik
	sepanjang $g$ secara bertahap sebagai
	% segment-id: o014.li2.chapter2.subobjects-and-isomorphism-theorems.g017
	\[\begin{tikzcd}
		W \arrow[d, "{f'}"'] \arrow[r] & \Image(g) \dtimes{Z} Y \arrow[r, "{g''}"] \arrow[d] & Y \arrow[d, "f"] \\
		X \arrow[twoheadrightarrow, r] & \Image(g) \arrow[hookrightarrow, r] \arrow[phantom, lu, "\Box" description] & Z ; \arrow[phantom, lu, "\Box" description]
	\end{tikzcd}\]
	Proposisi~\ref{prop:Abel-cat-pull-push} menyatakan bahwa
	$W\to\Image(g)\dtimes{Z}Y$ merupakan epimorfisme, sehingga
	$\Coker(g')=\Coker(g'')$. Dengan demikian, cukup ditinjau kasus
	ketika $f$ dan $g$ keduanya merupakan monomorfisme. Dalam hal ini
	$W=Y\cap X$, dan $c$ berfaktor sebagai
	$Y/(Y\cap X)\rightiso(Y+X)/X\hookrightarrow Z/X$.
\end{bukti}

% segment-id: o014.li2.chapter2.subobjects-and-isomorphism-theorems.th002
\begin{teorema}\label{prop:subobject-modularity}
	Untuk setiap objek $X$, himpunan terurut parsial
	$(\mathrm{Sub}_X,\subset)$ merupakan kisi modular terbatas dalam
	pengertian Definisi~\ref{def:modular-lattice}.
\end{teorema}
% segment-id: o014.li2.chapter2.subobjects-and-isomorphism-theorems.pf007
\begin{bukti}
	Kita sudah mengetahui bahwa $(\mathrm{Sub}_X,\subset)$ merupakan kisi
	terbatas. Sekarang tinjau sebarang interval $[Y_1,Y_2]$ dalam
	$\mathrm{Sub}_X$. Lema~\ref{prop:modularity-criterion} mereduksi persoalan
	menjadi pernyataan berikut: untuk setiap $A,B,C\in[Y_1,Y_2]$,
	% segment-id: o014.li2.chapter2.subobjects-and-isomorphism-theorems.q008
	\begin{equation}\label{eqn:Abel-cat-interval-complement}
		\left( A, B \;\text{keduanya merupakan komplemen dari}\; C, \; A \subset B \right) \implies A = B.
	\end{equation}

	Pertama-tama, bijeksi yang mempertahankan urutan dalam
	Teorema~\ref{prop:Abel-cat-isom-thm}(ii) mereduksi
	\eqref{eqn:Abel-cat-interval-complement} ke kasus $Y_1=0$ dan $Y_2=X$.
	Dengan Proposisi~\ref{prop:biproduct-sum-intersection}, hipotesis mengenai
	komplemen menjadi
	$A\oplus C\rightiso X\leftiso B\oplus C$, dengan kedua isomorfisme
	diinduksi oleh monomorfisme
	$\iota_A,\iota_B,\iota_C$ dari $A,B,C$ menuju $X$. Menurut hipotesis,
	terdapat $\alpha:A\to B$ sedemikian sehingga
	$\iota_A=\iota_B\alpha$. Maka terdapat diagram komutatif
	% segment-id: o014.li2.chapter2.subobjects-and-isomorphism-theorems.g018
	\[\begin{tikzcd}[column sep=large]
		B \oplus C \arrow[r, "{(\iota_B, \iota_C)}"] & X \\
		A \oplus C \arrow[ru, "{(\iota_A, \iota_C)}"'] \arrow[u, "{(\alpha, \identity_C)}"] &
	\end{tikzcd}\]
	Jadi, $(\alpha,\identity_C)$ merupakan isomorfisme. Menurut
	Lema~\ref{prop:isomorphism-matrix}, hal ini selanjutnya mengakibatkan bahwa
	$\alpha$ merupakan isomorfisme. Dengan demikian,
	\eqref{eqn:Abel-cat-interval-complement} terbukti.
\end{bukti}

% segment-id: o014.li2.chapter2.subobjects-and-isomorphism-theorems.p008
Sifat kisi modular merupakan jembatan penting untuk memindahkan sejumlah
argumen standar dari teori modul ke kategori abelian.
